\documentclass{article}
\usepackage{amsmath, amssymb, amsthm}
\usepackage{geometry}
\geometry{a4paper, margin=1in}
\usepackage{fancyhdr}
\usepackage{hyperref}
\usepackage[russian]{babel}
\usepackage[utf8]{inputenc}
\usepackage{listings}
\usepackage{xcolor}

\pagestyle{fancy}
\fancyhf{}
\fancyhead[L]{Algorithms}
\fancyhead[R]{\thepage}

\title{Algorithms 1 \\ Homework 6}
\author{Nikita Zagainov}
\date{\today}

\lstset{
    language=Python,
    basicstyle=\ttfamily\small,
    keywordstyle=\color{blue},
    stringstyle=\color{red},
    commentstyle=\color{gray},
    backgroundcolor=\color{lightgray!20},
    frame=single,
    rulecolor=\color{black},
    numbers=left,
    numberstyle=\tiny\color{gray},
    stepnumber=1,
    numbersep=10pt,
    showspaces=false,
    showstringspaces=false,
    breaklines=true,
    breakatwhitespace=true,
    tabsize=4,
    captionpos=b
}

\begin{document}

\maketitle

\section{Разные задачи. Расширенный алгоритм Евклида.}
\begin{enumerate}
    \item
          По формуле суммы первых $n$ членов
          геометрической прогрессии имеем:
          \[
              f(n) = \frac{c^{n + 1} - 1}{c - 1} = \Theta(c^n)
          \]

    \item
          Рассмотрим характеристическое уравнение,
          полученное из рекуррентного соотношения \\
          $T(n) = T(n - 1) + 4T(n - 3)$:
          \[
              r^n = r^{n - 1} + 4r^{n - 3} \Longleftrightarrow
              r^3 = r^2 + 4
          \]
          Корни характеристического уравнения:
          \[
              r_1 = 2,
              \quad r_2 = -\frac{1}{2} + \frac{i\sqrt{7}}{2},
              \quad r_3 = -\frac{1}{2} - \frac{i\sqrt{7}}{2}
          \]
          Тогда решение рекуррентного соотношения представимо
          в виде из линейной комбинации:
          \[
              T(n) = c_1 2^n +
              c_2 \left(-\frac{1}{2} + \frac{i\sqrt{7}}{2}\right)^n +
              c_3 \left(-\frac{1}{2} - \frac{i\sqrt{7}}{2}\right)^n
          \]
          Это экспоненциальная функция, а значит,
          взяв логарифм от обеих частей, справа
          получим линейную функцию:
          \[
              \log T(n) = O(n)
          \]

    \item
          Наименьшее общее кратное двух чисел $a$ и $b$ равно
          их произведению, делённому на их наибольший общий делитель:
          \[
              \mathrm{lcm}(a, b) = \frac{a \cdot b}{\gcd(a, b)}
          \]
          Если для нахождения НОД используется алгоритм Евклида,
          то алгоритм нахождения НОК можно записать следующим образом:
          \begin{lstlisting}
def gcd(a, b):
    if b == 0:
        return a
    return gcd(b, a % b)

def lcm(a, b):
    return a * b // gcd(max(a, b), min(a, b))
          \end{lstlisting}

          Оценим асимптотику алгоритма для $a$ и $b$ с
          количеством битов $n$ и $m$ соответственно.
          Асимптотически количество вызовов алгоритма
          Евклида равно
          \[
              O(\log \min(a, b)) = O(\min(n, m)) = O(m)
          \]
          так как на каждом шаге алгоритма Евклида одно из чисел
          уменьшается как минимум вдвое.
          Таким образом, асимптотика алгоритма алгоритма
          Евклида равна $O(nm^2)$, так как на каждом шаге
          алгоритма Евклида происходит деление с остатком
          за $O(nm)$.

          Тогда асимптотика алгоритма нахождения НОК равна
          $O(nm + nm^2 + m(n + m)) = O(nm^2)$, где произведение
          и требует $O(nm)$ времени,
          деление --- $O(m(n + m))$,
          сравнение --- линейное.

    \item
          Предположим, что при больших $n$ справедливо
          $T(n) \leq Cn$ для некоторой константы $C$.
          Тогда:
          \begin{align*}
               & T(n) = 3T(\frac{n}{4}) + T(\frac{n}{6}) + n \\
               & T(n) \leq 3C\frac{n}{4} + C\frac{n}{6} + n  \\
               & T(n) \leq Cn\frac{11}{12} + n
          \end{align*}
          Для $C \geq 12$ верно:
          \[
              T(n) \leq Cn\frac{11}{12} + n \leq Cn
          \]
          Тогда $T(n) = \Theta(n)$.

    \item
          Расширенный алгоритм Евклида восстанавливает
          коэффициенты $x$ и $y$ в линейном представлении
          НОД $a$ и $b$ по следующему правилу:
          \[
              x_t, \ y_t = y_{t+1}, \ x_{t+1} -
              \lfloor \frac{a}{b} \rfloor y_{t+1}
          \]
          Построим таблицу рекурсивных вызовов:
          \begin{center}
              \begin{tabular}{|c|c|c|c|c|c|c|}
                  \hline
                  $t$ & $a$ & $b$ & $d$ & $\lfloor \frac{a}{b} \rfloor$ & $x$ & $y$ \\
                  \hline
                  \hline
                  0   & 45  & 36  & 9   & 1                             & 1   & -1  \\
                  \hline
                  1   & 36  & 9   & 9   & 4                             & 0   & 1   \\
                  \hline
                  2   & 9   & 0   & 9   & -                             & 1   & 0   \\
                  \hline
              \end{tabular}
          \end{center}
          Тогда $x = -1, \ y = 1$ (В таблице $a$ и $b$
          поменялись местами для удобства):
          \[
              \gcd(36, 45) = -1 \cdot 36 + 1 \cdot 45
          \]

    \item
          Пусть оптимальный алгоритм перекладывает
          $n$ колец за $T(n)$ шагов. Тогда, чтобы переложить
          последнее кольцо со стартового стержня на целевой,
          необходимо сначала переложить $n - 1$ кольцо на
          вспомогательный стержень за $T(n - 1)$ шагов,
          потом за один шаг переложить последнее кольцо
          на целевой стержень, и затем переложить $n - 1$
          кольцо с вспомогательного стержня на целевой за
          $T(n - 1)$ шагов. Таким образом, рекуррентное
          соотношение:
          \[
              T(n) = 2T(n - 1) + 1
          \]

          Откуда $T(n) = 2^n - 1 = \Theta(2^n)$.
          Алгоритм, делающий $2^n - 1$ шагов:
          \begin{lstlisting}
from collections import deque


def place_rings(n: int, initial: deque, target: deque, helper: deque):
    if n == 1:
        target.append(initial.pop())
    else:
        place_rings(n - 1, initial, helper, target)
        target.append(initial.pop())
        place_rings(n - 1, helper, target, initial)
          \end{lstlisting}

    \item
          Для начала создадим копию входного массива
          и отсортируем его за $O(n \log n)$. Далее
          инициализируем указатель на первый элемент
          один раз пройдём другим указателем по длине
          массива, постоянно сравнивая подмассивы
          элементов из входного массива с начала
          и до указателя с подмассивы элементов из
          отсортированного массива от начала и до
          указателя. Если подмассивы равны, то подмассив
          от первого указателя до второго является
          одним из возможных подмассивов. Добавляем этот
          отрезок к ответу и записываем значение
          второго указателя в первый указатель.

          Этот алгоритм гарантированно найдёт максимальное
          количество подмассивов. Если это оказалось не так,
          то в ответе алгоритма найдётся подмассив, который
          может быть поделён на меньшие подмассивы, но
          это противоречит тому, что алгоритм записывает
          подмассивы в результат каждый раз, когда подмассивы
          элементов совпадают.

          Асимптотика алгоритма:
          $O(n \log n + n) = O(n \log n)$ на сортировку
          и один проход по массиву.

          \begin{lstlisting}
def max_sub_arrays(arr: list):
    sorted_arr = sorted(arr)
    sub_arrays = []

    arr_running_sum = 0
    sorted_arr_running_sum = 0
    current_start = 0

    for i in range(len(arr)):
        arr_running_sum += arr[i]
        sorted_arr_running_sum += sorted_arr[i]

        if arr_running_sum == sorted_arr_running_sum:
            sub_arrays.append([current_start, i + 1])
            current_start = i + 1
            arr_running_sum = 0
            sorted_arr_running_sum = 0

    return sub_arrays
          \end{lstlisting}

          Этот алгоритм сравнивает не подмассивы элементов,
          а суммы одинаковых префиксов массивов.
          Такой подход тоже работает, так как префиксная
          сумма отсортированного массива --- сумма такого
          же количества минимальных элементов в исходном
          массиве. Тогда суммы будут равно тогда и
          только тогда, когда подмассивы равны. Для всех
          отрезков элементов после первого элемента это
          свойство сохраняется, так как после нахождения
          первого подмассива процесс эквивалентен нахождению
          первого подмассива в оставшемся массиве.

\end{enumerate}
\end{document}